\section{État de l'Art}

\subsection{IA Générative dans l'Idéation et la Génération de Concepts}

\subsubsection{Modèles Text-to-Image}

Les modèles de génération d'images à partir de texte ont connu une évolution rapide ces dernières années. DALL-E \cite{ramesh2021zero}, Midjourney et Stable Diffusion \cite{rombach2022high} sont de plus en plus utilisés dans les premières étapes de conception pour l'inspiration, la création de mood boards et l'idéation conceptuelle \cite{bartlett2024generative, takale2024advancements}.

\textbf{Stable Diffusion}, en particulier, s'est imposé comme une solution open-source performante basée sur les modèles de diffusion latente (Latent Diffusion Models - LDM). Son architecture permet :

\begin{itemize}[leftmargin=*]
    \item Une génération d'images haute résolution (512x512 à 1024x1024)
    \item Un contrôle fin via le prompt engineering
    \item Une extensibilité avec des modules comme ControlNet
    \item Une exécution possible sur hardware grand public
\end{itemize}

\subsubsection{Prompt Engineering et Optimisation}

Kwon et al. \cite{kwon2024designer} proposent un processus d'idéation "Designer-Generative AI" qui met l'accent sur la formulation de prompts et l'exploration itérative d'images pour aligner les sorties de l'IA avec l'intention du designer. Leur travail souligne l'importance de structurer les prompts et d'utiliser l'IA pour l'extraction et le raffinement de mots-clés.

Cette approche est partiellement reflétée dans DesignPro AI à travers :
\begin{itemize}[leftmargin=*]
    \item Le module \texttt{DesignPromptGenerator} utilisant Mistral AI
    \item L'intégration de méthodologies de design dans la génération de prompts
    \item Un système de raffinement itératif basé sur le feedback utilisateur
\end{itemize}

\subsection{Collaboration Humain-IA dans le Design}

\subsubsection{Paradigme de Co-création}

Le consensus dans la littérature récente est que l'IA servira de partenaire puissant pour les designers humains, augmentant leurs capacités plutôt que de les remplacer \cite{bartlett2024generative, elmontassir2024leveraging}.

EIMT \cite{eimt2025future} décrit l'IA générative se transformant en "partenaire ou collaborateur", où artistes, écrivains et designers s'engagent avec l'IA pour repousser les frontières créatives.

\subsubsection{Frameworks d'Augmentation Humain-IA}

Des frameworks émergent pour évaluer et optimiser cette collaboration. Le Human-AI Augmentation Index \cite{friendsofeurope2025} se concentre sur la façon dont l'IA complète les capacités humaines, favorisant l'innovation en permettant aux humains de se concentrer sur des défis complexes et créatifs.

DesignPro AI, bien qu'il ne soit pas un système multi-agents, favorise un environnement collaboratif à travers :
\begin{itemize}[leftmargin=*]
    \item Des interfaces interactives permettant le guidage de l'IA
    \item Un système de scoring pour orienter les générations
    \item Des boucles de feedback itératives
    \item Une transparence sur les décisions de l'IA
\end{itemize}

\subsection{Outils de Design Génératif Existants}

\subsubsection{Solutions Commerciales}

Des outils commerciaux comme Autodesk Generative Design (dans Fusion 360) et Siemens NX exploitent déjà l'IA pour l'optimisation structurelle et l'exploration d'alternatives de design basées sur des contraintes d'ingénierie \cite{buonamici2020generative, rawat2023}.

Ces outils se concentrent généralement sur la génération de géométrie 3D et l'optimisation topologique. DesignPro AI, tout en visant également l'optimisation du design, se concentre principalement sur :
\begin{itemize}[leftmargin=*]
    \item La visualisation conceptuelle 2D en phase amont
    \item L'analyse textuelle DfX
    \item L'intégration de méthodologies de design
\end{itemize}

\subsubsection{Tendances vers la Démocratisation}

La tendance est vers des outils d'IA plus démocratisés \cite{eimt2025future, parivedaolutions2025}. L'utilisation par DesignPro AI de modèles open-source (Stable Diffusion, Mistral) et d'une stack web moderne (Next.js, Supabase) s'aligne avec cette vision, rendant potentiellement les capacités avancées plus accessibles.

Codewave \cite{codewave2025} fournit un guide sur l'utilisation de l'IA générative dans la conception de produits, soulignant son rôle dans :
\begin{itemize}[leftmargin=*]
    \item La génération de multiples variations
    \item La réduction du gaspillage de matériaux
    \item La personnalisation de masse
\end{itemize}

\subsection{Design for X (DfX) et IA}

\subsubsection{Principes DfX}

Les méthodologies Design for X regroupent un ensemble de pratiques visant à optimiser un produit selon différents critères \cite{kuo2001}:

\begin{description}[leftmargin=*]
    \item[DFM (Design for Manufacturing)] : Simplification de la fabrication, réduction des coûts de production
    \item[DFA (Design for Assembly)] : Minimisation du nombre de pièces, facilitation de l'assemblage
    \item[DFS (Design for Serviceability)] : Facilitation de la maintenance et des réparations
    \item[DFSust (Design for Sustainability)] : Réduction de l'impact environnemental, recyclabilité
\end{description}

\subsubsection{Intégration de DfX avec l'IA}

El Montassir et al. \cite{elmontassir2024leveraging} proposent un framework DfX intégré avec l'IA générative pour l'innovation et l'efficacité. Leur approche souligne l'importance d'intégrer les considérations DfX dès les phases précoces de conception.

DesignPro AI innove en :
\begin{itemize}[leftmargin=*]
    \item Intégrant explicitement les principes DfX dans la génération de prompts
    \item Fournissant une évaluation DfX automatisée via Agent Q
    \item Générant des rapports R.E.A.L. avec analyses FEA/DFM
    \item Permettant une sélection de méthodologie DfX par l'utilisateur
\end{itemize}

\subsection{Analyse Visuelle et Évaluation de Design}

\subsubsection{Vision par Ordinateur pour le Design}

L'analyse automatique de designs visuels est un domaine en développement. Les approches traditionnelles incluent :
\begin{itemize}[leftmargin=*]
    \item Détection de caractéristiques géométriques
    \item Analyse de complexité visuelle
    \item Évaluation esthétique computationnelle
\end{itemize}

\subsubsection{Modèles Vision-Language}

L'émergence de modèles vision-language comme CLIP \cite{radford2021learning} et GPT-4 Vision \cite{openai2023gpt4v} permet une analyse sémantique plus riche des images.

DesignPro AI exploite GPT-4 Vision pour :
\begin{itemize}[leftmargin=*]
    \item Analyser les formes et proportions
    \item Identifier les matériaux apparents
    \item Évaluer l'ergonomie visible
    \item Détecter les défauts ou incohérences visuelles
\end{itemize}

Cette approche représente une avancée par rapport aux méthodes purement contextuelles, atteignant une précision d'évaluation de ~90\% contre ~40\% pour l'analyse contextuelle seule.

\subsection{Lacunes Adressées par DesignPro AI}

L'analyse de l'état de l'art révèle plusieurs lacunes que DesignPro AI tente de combler :

\begin{enumerate}[leftmargin=*]
    \item \textbf{Intégration fragmentée} : Les outils existants se concentrent soit sur la génération visuelle, soit sur l'optimisation technique, rarement les deux
    
    \item \textbf{Absence de DfX précoce} : Les considérations de fabricabilité sont généralement introduites tardivement dans le processus
    
    \item \textbf{Manque d'analyse visuelle} : Les systèmes d'évaluation se basent souvent uniquement sur des métriques textuelles
    
    \item \textbf{Workflow non unifié} : Les designers doivent jongler entre multiples outils (ChatGPT pour prompts, Midjourney pour images, outils CAO pour validation)
    
    \item \textbf{Accessibilité limitée} : Les solutions commerciales sont coûteuses et propriétaires
\end{enumerate}

\subsection{Positionnement de DesignPro AI}

DesignPro AI se positionne comme une plateforme intégrée qui :

\begin{itemize}[leftmargin=*]
    \item \textbf{Combine} génération visuelle (Stable Diffusion) et analyse textuelle (Mistral AI)
    \item \textbf{Intègre} explicitement les méthodologies DfX dès la phase d'idéation
    \item \textbf{Fournit} une analyse visuelle réelle via GPT-4 Vision
    \item \textbf{Offre} un workflow complet en 4 phases dans une interface unique
    \item \textbf{Utilise} des technologies open-source pour une meilleure accessibilité
    \item \textbf{Supporte} l'itération et le raffinement via un système de scoring
\end{itemize}

Le tableau \ref{tab:comparison} compare DesignPro AI avec les approches existantes.

\begin{table}[H]
\centering
\caption{Comparaison de DesignPro AI avec les approches existantes}
\label{tab:comparison}
\begin{tabular}{@{}lcccc@{}}
\toprule
\textbf{Caractéristique} & \textbf{Midjourney} & \textbf{Autodesk} & \textbf{GENAI-DFX} & \textbf{DesignPro AI} \\ 
\midrule
Génération d'images & \checkmark & & \checkmark & \checkmark \\
Prompt engineering & Manuel & & \checkmark & \checkmark (Mistral) \\
Sketch-to-image & & & \checkmark & \checkmark (ControlNet) \\
Analyse DfX & & \checkmark & \checkmark & \checkmark (Agent Q) \\
Analyse visuelle & & & & \checkmark (GPT-4 Vision) \\
Simulation FEA/DFM & & \checkmark & & \checkmark (R.E.A.L.) \\
Workflow intégré & & & \checkmark & \checkmark \\
Open-source & & & \checkmark & \checkmark \\
Interface web & & & \checkmark & \checkmark (Next.js) \\
\bottomrule
\end{tabular}
\end{table}

\subsection{Défis et Directions Futures}

La littérature identifie plusieurs défis persistants \cite{torabi2024, addepto2024} :

\begin{itemize}[leftmargin=*]
    \item \textbf{Interprétabilité} : Comprendre comment l'IA génère ses sorties
    \item \textbf{Biais} : Préjugés intégrés dans les modèles d'IA
    \item \textbf{Propriété intellectuelle} : Questions de droits d'auteur sur les créations IA
    \item \textbf{Validation} : Métriques pour mesurer l'impact de l'IA sur la créativité
    \item \textbf{Éthique} : Utilisation responsable de l'IA dans le design
\end{itemize}

DesignPro AI aborde certains de ces défis en :
\begin{itemize}[leftmargin=*]
    \item Maintenant l'humain dans la boucle de décision
    \item Fournissant des explications sur les évaluations Agent Q
    \item Documentant les modèles et méthodologies utilisés
    \item Permettant un contrôle fin via les paramètres de génération
\end{itemize}
